% Copyright (c) 2026 by C. Klukas.
% Licensed under the MIT License. See LICENSE for details.

This appendix lists the five example \paglet{}s the paper refers to,
with commentary that points at the lines that matter: four of them
complete, pi in the sections the text discusses. The
listings are generated from the sources in \code{cpp/examples/} by
\code{paglets-docs-sync}: a whole file appears without its license
header, and the \code{-{}-8<-{}-} marker comments that delimit the
excerpts of the main text are left out. An excerpt that the main text
already shows is referenced rather than repeated. Every example has a page in the
documentation \cite{pagletsCppDocs} with its recorded test results; the
tests named there run the four mesh examples on a mesh of hosts inside
one test process, and the counter in the WebAssembly engine alone and
through the command line. The chronicle page also records a live run on
three machines (Section~\ref{sec:evaluation}).

\subsection{Building and running a \paglet{}}\label{sec:examples-build}

The build of a \paglet{} is one \code{paglets\_add\_module} call in
\code{cpp/examples/CMakeLists.txt}; Listing~\ref{lst:app-cmake} shows the
calls for the counter and for pi. \code{SOURCES} names the C++ file.
\code{SCHEMA\_HEADER} and \code{SCHEMA\_NAMESPACE} name the header with the
message types and its namespace; the build runs \code{schema\_gen} on it
and writes \code{NAME.schema.gen.hpp}, the header the \paglet{} includes,
with the MessagePack encoders and the schema description
\code{paglets\_schema\_json}. \code{SERVICES} names the system services
whose generated clients the \paglet{} uses (Section~\ref{sec:mesh}).
\code{PATTERNS} adds the six services the headers of the patterns
library use (\code{directory}, \code{files}, \code{grants},
\code{locator}, \code{mesh\_info} and \code{user\_info}); the headers
themselves are on every module's include path. The zombie pattern also
needs \code{zombies}, which the chronicle lists itself. \code{INTERFACE}
names the interface description (Section~\ref{sec:tools}), which the
build puts into the module as a custom section. The module ends up as
\code{NAME.wasm} in \path{build/<preset>/guests/}. Guests are compiled for
\code{wasm32-wasip1} with \code{-std=c++23}, \code{-fno-exceptions} and
\code{-fno-rtti}.
Outside the repository, a project installs paglets/cpp with
\code{cmake -{}-install}, calls \code{find\_package} for \code{paglets}
0.2 and uses the same function.

% hand-copied from cpp/examples/CMakeLists.txt lines 6-17
\begin{lstlisting}[style=paglets-cmake,caption={The build of the counter and of pi \src{cpp/examples/CMakeLists.txt}},label={lst:app-cmake}]
paglets_add_module(counter
    SOURCES counter/counter.cpp
    SCHEMA_HEADER counter/messages.hpp
    SCHEMA_NAMESPACE counter_msgs
    INTERFACE counter/interface.json)
paglets_add_module(pi
    SOURCES pi/pi.cpp
    SCHEMA_HEADER pi/pi_msgs.hpp
    SCHEMA_NAMESPACE pi_msgs
    SERVICES mesh_info compute_slots
    PATTERNS
    INTERFACE pi/interface.json)
\end{lstlisting}

The counter runs on one host without a mesh (Listing~\ref{lst:app-run}).
Message bodies on the command line are JSON; the host passes them to the
\paglet{} as MessagePack and prints replies as JSON, enums as their names.
With \code{-{}-state-dir} and \code{-{}-keep} the \paglet{} outlives the
host process, and \code{call} resumes it from its last memory image in a
new process. The other four examples need a mesh of hosts that serve
(Section~\ref{sec:mesh}). They are launched on one host with
\code{remote launch}, which prints the \paglet{}'s ID, and driven with
\code{remote call}, which reaches a \paglet{} on the host the command is
connected to only: while the journey is away, its home host answers
\code{not\_found}, and the command is repeated until it is back
(messages between \paglet{}s follow a move,
Section~\ref{sec:communication}). The inventory and pi also need a
policy rule that allows the owner's \paglet{}s the operations they use
(\code{server-info}, and \code{files} for pi's output).

% the commands of docs/examples/counter.md and docs/examples/mesh-journey.md
\begin{lstlisting}[style=paglets-bash,caption={Running the counter locally, and the journey on a mesh},label={lst:app-run}]
H=build/macos-arm64/host/paglets-host
$H run build/macos-arm64/guests/counter.wasm \
    --call increment '{"by": 5, "note": "first"}' \
    --call increment '{"by": 3, "note": "triple", "mode": "multiply"}' \
    --call status
# The value is 5, then 15; status reports history_size 2 and average_step 4.

$H run build/macos-arm64/guests/counter.wasm --state-dir state --keep \
    --call increment '{"by": 5, "note": "a"}' --call increment '{"by": 2, "note": "b"}'
$H call --state-dir state all increment '{"by": 1, "note": "c"}'   # value 8

R="--connect https://lab-1:7443 --key olga.key --ledger ledger"
ID=$($H remote launch $R build/macos-arm64/guests/journey.wasm)
$H remote call $R $ID start '{"hosts": [], "label": ""}'
$H remote call $R $ID journal      # not_found while it is away; the journal once it is home
\end{lstlisting}

\subsection{Counter}\label{sec:examples-counter}

The counter (Listing~\ref{lst:app-counter}) keeps a value and the
history of the steps that changed it. It shows the smallest complete
setup: state that is ordinary C++ objects, a schema header
(Listing~\ref{lst:counter-messages} in Section~\ref{sec:model-schemas}),
optional start arguments, and an interface description
(Listing~\ref{lst:counter-interface} in
Section~\ref{sec:tools-interfaces}).
It is the \paglet{} of the image, state-directory and move tests,
because one \code{status} reply tells whether the state came through a
save, a move or a restart.

% snippet: cpp/examples/counter/counter.cpp [caption={The counter},label={lst:app-counter}]
\begin{lstlisting}[style=paglets-cpp,caption={The counter \src{cpp/examples/counter/counter.cpp}},label={lst:app-counter}]
// Counter example paglet. Its state is ordinary C++ objects on the heap and
// in static storage; memory images move it without any serialization code.

#include "counter.schema.gen.hpp"

#include <paglets/paglet.hpp>

#include <string>
#include <vector>

namespace {

using namespace counter_msgs;

class Counter : public paglets::Paglet {
public:
    Counter() {
        router().on<Increment>("increment", [this](const Increment& inc, paglets::Message& m) {
            value_ = inc.mode == Mode::add ? value_ + inc.by : value_ * inc.by;
            history_.push_back(inc);
            m.reply(status());
        });
        router().on("status", [this](paglets::Message& m) { m.reply(status()); });
        router().on<Echo>("echo", [](const Echo& echo, paglets::Message& m) { m.reply(echo); });
        router().on<Bloat>("bloat", [this](const Bloat& bloat, paglets::Message& m) {
            ballast_.emplace_back(std::size_t{bloat.kilobytes} * 1024, bloat.fill);
            m.reply(status());
        });
        router().on("schema", [](paglets::Message& m) { m.reply(std::string(paglets_schema_json)); });
        router().on<std::string>("log", [this](const std::string& text, paglets::Message& m) {
            paglets::log(text);
            m.reply(status());
        });
    }

    // Optional start arguments (`--args '{"value": 10}'`).
    void on_created(paglets::Started& s) override {
        Start start;
        if (!s.args.empty() && s.decode_args(start)) value_ = start.value;
    }

private:
    Status status() const {
        Status s;
        s.value = value_;
        s.history_size = static_cast<std::uint32_t>(history_.size());
        const std::size_t first = history_.size() > 3 ? history_.size() - 3 : 0;
        for (std::size_t i = first; i < history_.size(); ++i) {
            s.recent_notes.push_back(history_[i].note);
        }
        if (!history_.empty()) {
            double sum = 0;
            for (const auto& h : history_) sum += static_cast<double>(h.by);
            s.average_step = sum / static_cast<double>(history_.size());
        }
        s.memory_pages = static_cast<std::uint32_t>(__builtin_wasm_memory_size(0));
        return s;
    }

    std::int64_t value_ = 0;
    std::vector<Increment> history_;
    std::vector<std::vector<std::uint8_t>> ballast_;
};

}  // namespace

PAGLETS_PAGLET(Counter)
\end{lstlisting}
% /snippet

The constructor registers six handlers. The typed form
\code{router().on<Increment>} decodes the body into the struct before the
lambda runs and answers \code{malformed} when it cannot
(Section~\ref{sec:model}); the untyped form of \code{status} gets only the
\code{Message}. \code{on\_created} runs once, at creation, and decodes
the start arguments if there are any (Section~\ref{sec:lifecycle}). The
three members at the end of the class are the whole state: an integer, a
vector of structs with strings in them, and a vector of byte vectors that
\code{bloat} fills for experiments with the image size. Nothing in the
file is about serialization or movement. \code{status()} reports
\code{\_\_builtin\_wasm\_memory\_size(0)}, the size of the linear memory
in pages of 64\,KB, which is almost all of the image a move carries; the
rest is the exported mutable globals, here one, the stack pointer. A move
does not carry the whole image: the destination fetches only the pages
it lacks, each compressed (Section~\ref{sec:movement}).
Section~\ref{sec:evaluation} records the image sizes.

The schema header (Listing~\ref{lst:counter-messages}) is plain C++: an
enum, and structs with integers of
several widths, floating-point numbers, strings, byte vectors, vectors of
structs and \code{std::optional} fields. \code{Echo} has a field of
every supported kind of type and comes back unchanged, which is how the
tests check the generated encoder of the guest against the
reflection-based one of the host. The interface description
(Listing~\ref{lst:counter-interface}) names the
start type \code{Start}, marks \code{increment} and \code{log} as
\code{both} (a one-way send or a call) and the others as \code{call},
and defines three records itself, which is how their fields get a
\code{doc} and a \code{default}. The build merges the other records and
the enum of the schema header (\code{Echo}, \code{Status}, \code{Mode})
into \code{types} when it puts the description into the module, a type
the file defines taking precedence, so \code{paglets-web}
(Section~\ref{sec:tools}) has a form for every message.

\subsection{Mesh Journey}\label{sec:examples-journey}

The journey (Listing~\ref{lst:app-journey}; its messages and journal are
Listing~\ref{lst:app-journey-messages}) is the first \paglet{} to read
for movement. It visits every
host of the mesh in turn, or the hosts it is asked for, writes down each
stop in a journal, and comes home. The journal is a \code{std::vector} of
structs in a member and travels with the \paglet{} inside the memory
image. The example also shows a deferred reply, the \code{mesh-info}
service, and two of the three hooks of a move, \code{on\_arrived} and
\code{on\_move\_failed} (\code{on\_dispatching} it does not need).

% snippet: cpp/examples/journey/journey.cpp [caption={The Mesh Journey},label={lst:app-journey}]
\begin{lstlisting}[style=paglets-cpp,caption={The Mesh Journey \src{cpp/examples/journey/journey.cpp}},label={lst:app-journey}]
// Mesh Journey (planning/cpp-demo-paglets.md, demo 1): the "hello world" of
// paglets/cpp. The paglet visits every host of the mesh in turn and keeps a
// journal of where it was in an ordinary std::vector; memory images carry
// it from host to host without serialization code. At the end it comes
// home; the journal shows every host that can receive and run paglets.

#include "journey.schema.gen.hpp"

#include <paglets/patterns/services.hpp>

#include <algorithm>
#include <string>
#include <vector>

namespace {

using namespace journey_msgs;
namespace svc = paglets::services;
namespace pt = paglets::patterns;
using paglets::Message;

class Journey : public paglets::Paglet {
public:
    Journey() {
        router().on<Start>("start", [this](const Start& s, Message& m) { start(s, m); });
        router().on("journal", [this](Message& m) { (void)m.reply(journal_); });
    }

    void on_arrived(const paglets::abi::ArrivedEvent&) override { record_and_go_on(); }

    void on_move_failed(const paglets::abi::MoveFailed& f) override {
        if (journal_.state != "travelling") return;
        journal_.skipped.push_back(f.destination + ": " + f.reason);
        // Home out of reach: the journey ends where it is.
        if (f.destination == home_) {
            journal_.state = "failed";
            return;
        }
        go_on();  // the next host
    }

private:
    static std::int64_t now_ms() { return paglets::now_ns(paglets::abi::Clock::wall) / 1'000'000; }

    void start(const Start& s, Message& m) {
        if (journal_.state == "travelling") {
            (void)m.reply(Started{false, "already travelling", itinerary_});
            return;
        }
        journal_ = Journal{"travelling", {}, {}};
        next_ = 0;
        auto reply = std::make_shared<paglets::Capability>(m.defer_reply());
        // The itinerary: the hosts asked for, else every host mesh-info knows.
        auto mi = paglets::service("mesh-info");
        if (!mi) {
            journal_.state = "failed";
            (void)paglets::reply_to(*reply, Started{false, "no mesh-info", {}});
            return;
        }
        svc::mesh_info::Client client{*mi};
        auto sent = client.landscape(
            svc::mesh_info::LandscapeRequest{},
            [this, s, reply](paglets::Result<svc::mesh_info::Landscape> l, Message&) {
                if (!l || l->hosts.empty()) {
                    journal_.state = "failed";
                    (void)paglets::reply_to(*reply, Started{false, "no landscape", {}});
                    return;
                }
                home_ = l->hosts.front().host;  // this host comes first
                itinerary_.clear();
                for (const auto& h : l->hosts) {
                    if (h.host == home_) continue;
                    const bool asked = s.hosts.empty() || std::ranges::find(s.hosts, h.host_name) != s.hosts.end() ||
                                       std::ranges::find(s.hosts, h.host) != s.hosts.end();
                    const bool labelled =
                        s.label.empty() || std::ranges::find(h.labels, s.label) != h.labels.end();
                    if (asked && labelled) itinerary_.push_back(h.host);
                }
                // Hosts asked for by name that mesh-info does not know are skipped.
                for (const auto& want : s.hosts) {
                    const bool known = std::ranges::any_of(l->hosts, [&](const auto& h) {
                        return h.host_name == want || h.host == want;
                    });
                    if (!known) journal_.skipped.push_back(want + ": unknown to mesh-info");
                }
                (void)paglets::reply_to(*reply, Started{true, {}, itinerary_});
                record_and_go_on();  // home is the first stop
            });
        if (!sent) {
            journal_.state = "failed";
            (void)paglets::reply_to(*reply, Started{false, pt::error_text("mesh-info", sent.error()), {}});
        }
    }

    void record_and_go_on() {
        if (journal_.state != "travelling") return;
        const std::int64_t arrived = now_ms();
        pt::here([this, arrived](paglets::Result<svc::mesh_info::Snapshot> s) {
            Stop stop;
            stop.arrived_ms = arrived;
            stop.transfer_ms = left_ms_ > 0 ? std::max<std::int64_t>(0, arrived - left_ms_) : 0;
            if (s) {
                here_ = s->host;
                stop.host = s->host;
                stop.host_name = pt::host_label(*s);
                stop.os = s->os;
                stop.architecture = s->architecture;
                stop.load_per_cpu = s->load_per_cpu;
            }
            journal_.stops.push_back(std::move(stop));
            if (!journal_.stops.empty() && journal_.stops.back().host == home_ && journal_.stops.size() > 1) {
                journal_.state = "home";
                return;
            }
            go_on();
        });
    }

    void go_on() {
        if (next_ >= itinerary_.size() && here_ == home_) {
            journal_.state = "home";  // nothing (more) to visit
            return;
        }
        std::string destination = next_ < itinerary_.size() ? itinerary_[next_++] : home_;
        left_ms_ = now_ms();
        if (auto moved = paglets::dispatch(destination); !moved) {
            if (destination == home_) {
                journal_.state = "failed";
                journal_.skipped.push_back("home: " + pt::error_text("dispatch", moved.error()));
                return;
            }
            journal_.skipped.push_back(destination + ": " + pt::error_text("dispatch", moved.error()));
            go_on();
        }
    }

    Journal journal_{"idle", {}, {}};
    std::vector<std::string> itinerary_;
    std::size_t next_ = 0;
    std::string home_;
    std::string here_;
    std::int64_t left_ms_ = 0;
};

}  // namespace

PAGLETS_PAGLET(Journey)
\end{lstlisting}
% /snippet

% snippet: cpp/examples/journey/journey_msgs.hpp [caption={The journey's messages and journal},label={lst:app-journey-messages}]
\begin{lstlisting}[style=paglets-cpp,caption={The journey's messages and journal \src{cpp/examples/journey/journey\_msgs.hpp}},label={lst:app-journey-messages}]
// Messages of the Mesh Journey demo (planning/cpp-demo-paglets.md, demo 1).

#pragma once

#include <cstdint>
#include <string>
#include <vector>

namespace journey_msgs {

// `start`: visit `hosts` in turn (host names or key IDs; empty: every host
// mesh-info knows, or those with `label`), then come home.
struct Start {
    std::vector<std::string> hosts;
    std::string label;
};

struct Started {
    bool accepted = false;
    std::string reason;
    std::vector<std::string> itinerary;  // key IDs, in order
};

// One stop of the journey.
struct Stop {
    std::string host;  // key ID
    std::string host_name;
    std::string os;
    std::string architecture;
    double load_per_cpu = 0;
    std::int64_t arrived_ms = 0;   // Unix milliseconds, this host's clock
    std::int64_t transfer_ms = 0;  // from leaving the previous host (clocks may differ)
};

struct Journal {
    std::string state;  // idle, travelling, home, failed
    std::vector<Stop> stops;
    std::vector<std::string> skipped;  // hosts it could not reach, with the reason
};

}  // namespace journey_msgs
\end{lstlisting}
% /snippet

\code{start} cannot answer at once, because the itinerary comes from
\code{mesh-info}. It takes the reply capability out of the message with
\code{m.defer\_reply()}, keeps it in a \code{shared\_ptr} that the
continuation captures, and answers through \code{reply\_to} when the
landscape is there. The host the \paglet{} runs on comes first in the
landscape; that host is home. From then on the \paglet{} is driven by
events, never by a handler that waits: \code{record\_and\_go\_on} asks
\code{mesh-info} which host it is on (\code{pt::here}), records a
\code{Stop} in the continuation, and calls \code{go\_on}, which calls
\code{paglets::dispatch}. The move takes effect when the handler returns
(Section~\ref{sec:movement}); a dispatch that fails at once returns an
error code, and the journey skips the host in the same handler. A move
that fails after the handler returned arrives as \code{on\_move\_failed}
with the destination and the reason. On the new host, \code{on\_arrived}
runs first and continues with the members as they were: \code{journal\_},
\code{itinerary\_}, \code{next\_}, \code{home\_} and \code{left\_ms\_} are
all in the image. The journal's \code{state}
goes from \code{idle} to \code{travelling} to \code{home}, or to
\code{failed} when \code{mesh-info} gives no landscape or home is out of
reach; \code{journal} can be asked at any time from the host the
\paglet{} is on (Listing~\ref{lst:app-run}).

\subsection{Inventory Collector}\label{sec:examples-inventory}

The inventory (Listings~\ref{lst:app-inventory}
and~\ref{lst:app-inventory-messages}) collects every host's system
summary, load, volumes and
largest processes, or finds the processes with a given name on every
host. It is the fan-out example: the \paglet{} clones itself to every
host, each clone asks the \code{server-info} system \paglet{} of its host
and reports to the parent, and the parent keeps the report. One module
has both roles.

% snippet: cpp/examples/inventory/inventory.cpp [caption={The Inventory Collector},label={lst:app-inventory}]
\begin{lstlisting}[style=paglets-cpp,caption={The Inventory Collector \src{cpp/examples/inventory/inventory.cpp}},label={lst:app-inventory}]
// Inventory Collector, Process Finder and Mesh Top
// (planning/cpp-demo-paglets.md, demos 8, 10 and 11): a clone on every
// host asks `server-info` (the mesh policy decides which operations) for
// the system summary, load, volumes and processes, in the same schema on
// macOS, Linux and Windows, and reports; the parent keeps the mesh's
// inventory. With a process name it finds those processes everywhere.

#include "inventory.schema.gen.hpp"

#include <paglets/patterns/fanout.hpp>
#include <paglets/services/server_info.gen.hpp>

#include <algorithm>
#include <memory>
#include <string>
#include <vector>

namespace {

using namespace inventory_msgs;
namespace svc = paglets::services;
namespace si = paglets::services::server_info;
namespace pt = paglets::patterns;
using paglets::Message;

class Inventory : public paglets::Paglet {
public:
    Inventory() {
        router().on<Collect>("start", [this](const Collect& c, Message& m) { start(c, m); });
        router().on("report", [this](Message& m) { (void)m.reply(report_); });
        fan_.attach(router(), [this](const pt::FanOut::Child& c) {
            HostInventory h;
            if (!c.ok || !paglets::decode(c.result, h)) {
                h = HostInventory{};
                h.error = c.error.empty() ? "no report" : c.error;
            }
            h.host = c.host;
            h.host_name = c.host_name.empty() ? c.destination : c.host_name;
            report_.hosts.push_back(std::move(h));
            report_.finished = fan_.finished();
            if (report_.finished) {
                std::ranges::sort(report_.hosts, [](const auto& a, const auto& b) { return a.host_name < b.host_name; });
            }
        });
    }

    void on_cloned(paglets::Started& s) override {
        Collect q;
        if (s.caps.empty() || !s.decode_args(q)) {
            (void)paglets::dispose();
            return;
        }
        parent_ = paglets::Endpoint(s.caps[0].release());
        query_ = q;
        pt::lookup("server-info", [this](paglets::Result<paglets::Endpoint> ep) {
            if (!ep) return finish(pt::error_text("server-info", ep.error()));
            info_ = *ep;
            summary();
        });
    }

    void on_move_failed(const paglets::abi::MoveFailed& f) override {
        if (!f.clone.empty()) fan_.clone_failed(f);
    }

private:
    void start(const Collect& c, Message& m) {
        report_ = Report{};
        auto reply = std::make_shared<paglets::Capability>(m.defer_reply());
        auto go = [this, c, reply](std::vector<std::string> hosts) {
            const auto sent = fan_.start(hosts, paglets::encode(c), c.timeout_ms);
            (void)paglets::reply_to(*reply, Started{sent > 0, sent > 0 ? "" : "no clone went out",
                                                    static_cast<std::int64_t>(hosts.size())});
        };
        if (!c.hosts.empty()) return go(c.hosts);
        auto mi = paglets::service("mesh-info");
        if (!mi) return go({""});
        svc::mesh_info::Client client{*mi};
        auto sent = client.landscape({}, [go](paglets::Result<svc::mesh_info::Landscape> l, Message&) {
            std::vector<std::string> hosts{""};
            if (l) {
                for (std::size_t i = 1; i < l->hosts.size(); ++i) hosts.push_back(l->hosts[i].host);
            }
            go(std::move(hosts));
        });
        if (!sent) go({""});
    }

    // -- a clone: summary, load, volumes, processes, then the report -------------------

    void summary() {
        si::Client c{info_};
        auto sent = c.summary({}, [this](paglets::Result<si::Summary> s, Message&) {
            if (!s) return finish(pt::error_text("summary", s.error()));
            inv_.os = s->os;
            inv_.os_version = s->os_version;
            inv_.architecture = s->architecture;
            inv_.cpu_model = s->cpu_model;
            inv_.cpus = s->cpu_count;
            inv_.memory_total = static_cast<std::int64_t>(s->memory_total);
            inv_.uptime_s = static_cast<std::int64_t>(s->uptime_s);
            load();
        });
        if (!sent) finish(pt::error_text("summary", sent.error()));
    }

    void load() {
        si::Client c{info_};
        auto sent = c.load({}, [this](paglets::Result<si::Load> l, Message&) {
            if (l) {
                inv_.cpu_percent = l->cpu_percent;
                inv_.memory_available = static_cast<std::int64_t>(l->memory_available);
            }
            volumes();
        });
        if (!sent) volumes();
    }

    void volumes() {
        si::Client c{info_};
        auto sent = c.volumes({}, [this](paglets::Result<si::Volumes> v, Message&) {
            if (v) {
                for (const auto& x : v->volumes) {
                    inv_.volumes.push_back(VolumeInfo{x.mount_point, x.filesystem, static_cast<std::int64_t>(x.total),
                                                      static_cast<std::int64_t>(x.available)});
                }
            }
            processes();
        });
        if (!sent) processes();
    }

    void processes() {
        si::Client c{info_};
        si::ProcessesRequest q;
        q.limit = query_.process.empty() ? static_cast<std::uint32_t>(std::max<std::int64_t>(0, query_.top)) : 10'000;
        auto sent = c.processes(q, [this](paglets::Result<si::Processes> p, Message&) {
            if (p) {
                inv_.processes_total = p->total;
                for (const auto& x : p->processes) {
                    if (!query_.process.empty() && x.name.find(query_.process) == std::string::npos) continue;
                    inv_.processes.push_back(ProcessInfo{x.pid, x.name, static_cast<std::int64_t>(x.memory)});
                }
            }
            finish({});
        });
        if (!sent) finish({});
    }

    void finish(std::string error) {
        if (!error.empty()) return pt::report_error(parent_, std::move(error));
        pt::report(parent_, paglets::encode(inv_));
        (void)paglets::after_raw(2000, "done", {});
        router().on("done", [](Message&) { (void)paglets::dispose(); });
    }

    // The parent.
    pt::FanOut fan_;
    Report report_;
    // A clone.
    paglets::Endpoint parent_;
    paglets::Endpoint info_;
    Collect query_;
    HostInventory inv_;
};

}  // namespace

PAGLETS_PAGLET(Inventory)
\end{lstlisting}
% /snippet

% snippet: cpp/examples/inventory/inventory_msgs.hpp [caption={The inventory's messages and report},label={lst:app-inventory-messages}]
\begin{lstlisting}[style=paglets-cpp,caption={The inventory's messages and report \src{cpp/examples/inventory/inventory\_msgs.hpp}},label={lst:app-inventory-messages}]
// Messages of the Inventory Collector, Process Finder and Mesh Top demos
// (planning/cpp-demo-paglets.md, demos 8, 10 and 11).

#pragma once

#include <cstdint>
#include <string>
#include <vector>

namespace inventory_msgs {

// `start`: collect from every host (or `hosts`): the system summary, load,
// volumes and the `top` largest processes; with `process`, only processes
// whose name contains it (Process Finder).
struct Collect {
    std::vector<std::string> hosts;
    std::string process;
    std::int64_t top = 5;
    std::int64_t timeout_ms = 20'000;
};

struct Started {
    bool accepted = false;
    std::string reason;
    std::int64_t hosts = 0;
};

struct VolumeInfo {
    std::string mount_point;
    std::string filesystem;
    std::int64_t total = 0;
    std::int64_t available = 0;
};

struct ProcessInfo {
    std::int64_t pid = 0;
    std::string name;
    std::int64_t memory = 0;
};

// One host, in the same shape on every operating system.
struct HostInventory {
    std::string host;
    std::string host_name;
    std::string error;
    std::string os;
    std::string os_version;
    std::string architecture;
    std::string cpu_model;
    std::int64_t cpus = 0;
    std::int64_t memory_total = 0;
    std::int64_t memory_available = 0;
    double cpu_percent = 0;
    std::int64_t uptime_s = 0;
    std::vector<VolumeInfo> volumes;
    std::int64_t processes_total = 0;
    std::vector<ProcessInfo> processes;
};

struct Report {
    bool finished = false;
    std::vector<HostInventory> hosts;
};

}  // namespace inventory_msgs
\end{lstlisting}
% /snippet

Which role an instance plays is decided by the hook that runs:
\code{on\_cloned} runs in a clone only. The clone's \code{Started}
carries the arguments (the \code{Collect} request, encoded by the parent)
and one capability, the parent's endpoint, which the clone takes with
\code{s.caps[0].release()}; a clone with a destination can carry only
endpoints to \paglet{}s and artifact capabilities
(Section~\ref{sec:movement-clones}). The parent's side is the
\code{FanOut} pattern: \code{fan\_.attach} registers the handlers for
\code{fanout.report} and \code{fanout.timeout} on the router, and
\code{fan\_.start} clones the \paglet{} to each destination with
\code{clone\_raw(args, \{self()\}, destination)}, drops its own handle to
the clone, and starts the deadline timer. \code{start} replies as soon
as the clones went out, with the number of hosts asked and whether at
least one clone went out; \code{report} returns the report at any time, with \code{finished} set once every clone reported
or timed out. A clone that never arrives comes back as
\code{on\_move\_failed} with the clone's ID, which
\code{fan\_.clone\_failed} turns into a finished child with an error.
Without a list of hosts, the parent takes every host of the landscape but
the first, its own, addressed as the empty string.

A clone looks \code{server-info} up through the directory
(\code{pt::lookup}); a policy rule allows or denies each of its
operations separately. The four requests follow each other, each started in
the continuation of the previous reply: the summary is required, the
others leave their fields empty when denied. \code{finish} reports with
\code{pt::report}, which first asks \code{mesh-info} where the clone is,
so that the parent can name the host, and then sends
\code{fanout.report} through the parent's endpoint. The clone then
disposes of itself two seconds later, through a timer and a handler
registered at run time. The timer is not needed for the report to leave:
a send is handed to the host inside the handler, and a \code{dispose}
takes effect only when the handler returns, so the \code{sent} callback
of \code{pt::report} would end the clone just as safely. Reports arrive
at the parent in any order; the parent sorts the hosts by name when the
last one is in.

\subsection{Pi}\label{sec:examples-pi}

Pi computes hexadecimal digits of $\pi$ on the hosts of the mesh with the
Bailey--Borwein--Plouffe formula, which gives the digits at any position
without the digits before them. The job \paglet{} stays where it was
started, cuts the digits into chunks, and sends each chunk in a clone of
itself to a host that \code{mesh-info} selects for at least one free
compute slot and a load of at most 4 per CPU, or its own host when none
fits. The clone, a worker, asks its host's \code{compute-slots} for a
slot, computes, sends its digits to the job and ends. The job takes the
chunks in order, writes them to files on its host when it was started
with an output root (otherwise the digits stay in its memory), sends a
chunk out again when no answer came in time, and at the end, unless told
otherwise, converts the digits into decimal ones. The file is 538 lines;
the sections below are the ones the text refers to (two of them appear
in the main text: the worker's slot request,
Listing~\ref{lst:pi-worker}, and the decimal conversion,
Listing~\ref{lst:pi-convert}), and the header with the messages follows
them whole (Listing~\ref{lst:app-pi-messages}).
As in the inventory, one module has both roles: \code{on\_cloned} sets
\code{worker\_}.

The job sends \code{n} chunks round robin over the hosts
\code{mesh-info} selected, each in a clone that gets its \code{Task} as
arguments and the job's endpoint as its only capability; the empty host
name is the job's own host, and a clone for it is made locally
(Listing~\ref{lst:app-pi-clone}):

% snippet: cpp/examples/pi/pi.cpp#clone [caption={Pi: a clone per chunk},label={lst:app-pi-clone}]
\begin{lstlisting}[style=paglets-cpp,caption={Pi: a clone per chunk \src{cpp/examples/pi/pi.cpp}},label={lst:app-pi-clone}]
// `n` chunks, round robin over `hosts` ("" is this host).
void send_to(const std::vector<std::string>& hosts, std::int64_t n) {
    std::size_t next = 0;
    for (auto& c : chunks_) {
        if (n <= 0) break;
        if (c.sent || c.done) continue;
        const std::string& host = hosts[next++ % hosts.size()];
        Task t{c.start, c.count, ++c.attempt, params_.work_ms};
        std::optional<std::string> destination;
        if (!host.empty() && host != host_) destination = host;
        auto worker = paglets::clone_raw(paglets::encode(t), {paglets::self()}, destination);
        if (!worker) {
            paglets::log("pi: no worker for chunk " + std::to_string(c.start) + ": " +
                         std::string(abi::error_name(worker.error())));
            break;
        }
        (void)worker->drop();  // the worker answers through its own endpoint to the job
        c.sent = true;
        c.sent_ms = now_ms();
        --n;
    }
}
\end{lstlisting}
% /snippet

A worker asks its host for one core and acts on the verdict
(Listing~\ref{lst:pi-worker} in Section~\ref{sec:mesh-compute}): it runs,
waits for \code{compute.granted}, or moves to the host the redirect names
and asks there again from \code{on\_arrived}; a host without
\code{compute-slots}, or a rejected request, lets it compute at once.

A handler may run for at most five seconds, the host's handler budget; a
handler that runs longer is stopped and its \paglet{} ends. So the
worker computes in slices of about 200\,ms and continues in the next
message, a timer to itself with no delay
(Listing~\ref{lst:app-pi-slices}):

% snippet: cpp/examples/pi/pi.cpp#slices [caption={Pi: computing in slices between messages},label={lst:app-pi-slices}]
\begin{lstlisting}[style=paglets-cpp,caption={Pi: computing in slices between messages \src{cpp/examples/pi/pi.cpp}},label={lst:app-pi-slices}]
// A slice of about 200 ms: a handler call may take at most 5 s, and far
// out a single block of six digits takes seconds. The rest follows in
// the next message.
void compute_more() {
    const std::int64_t began = now_ms();
    while (static_cast<std::int64_t>(digits_.size()) < task_.count) {
        if (block_.step(20'000)) {
            digits_ += block_.digits();
            const std::int64_t at = task_.start + static_cast<std::int64_t>(digits_.size());
            const std::int64_t left = task_.start + task_.count - at;
            if (left > 0) block_ = pi_bbp::Block(at, static_cast<int>(std::min<std::int64_t>(6, left)));
        }
        if (now_ms() - began >= 200 && static_cast<std::int64_t>(digits_.size()) < task_.count) {
            (void)paglets::after_raw(0, "compute-more", {});
            return;
        }
    }
    report_digits();
}
\end{lstlisting}
% /snippet

The job keeps a chunk's digits only until the chunk is next in order;
then they go into the file and into the fraction, and leave the job's
memory, because a \paglet{}'s memory holds at most 16\,MB and a message
at most 1\,MB (Listing~\ref{lst:app-pi-in-order}):

% snippet: cpp/examples/pi/pi.cpp#in_order [caption={Pi: chunks leave the job's memory in order},label={lst:app-pi-in-order}]
\begin{lstlisting}[style=paglets-cpp,caption={Pi: chunks leave the job's memory in order \src{cpp/examples/pi/pi.cpp}},label={lst:app-pi-in-order}]
// Chunks that are next in order leave the job's memory: into the file,
// and into the fraction for the decimal conversion.
void take_in_order() {
    while (next_ < chunks_.size() && chunks_[next_].done) {
        Chunk& c = chunks_[next_++];
        done_ += c.count;
        if (shown_hex_.size() < shown_digits) shown_hex_ += c.digits.substr(0, shown_digits - shown_hex_.size());
        if (params_.decimal) {
            for (char h : c.digits) add_nibble(static_cast<std::uint32_t>(h <= '9' ? h - '0' : h - 'A' + 10));
        }
        if (dir_) {
            hex_file_.pending += c.digits;
            flush(hex_file_);
        }
        std::string().swap(c.digits);
    }
    if (finished() && state_ == "computing") {
        if (params_.decimal) return start_conversion();
        state_ = "writing";
        maybe_done();
    }
}

// The fraction, 8 hex digits per 32-bit limb, the most significant first.
void add_nibble(std::uint32_t v) {
    partial_ = (partial_ << 4) | v;
    if (++partial_n_ == 8) {
        limbs_.push_back(partial_);
        partial_ = 0;
        partial_n_ = 0;
    }
}
\end{lstlisting}
% /snippet

Every 500\,ms the job checks the chunks in flight; one without an answer
for \code{chunk\_timeout\_ms} counts as waiting again and goes out with
the next batch, to the hosts \code{mesh-info} selects then (a host that is
gone drops out of the selection once its last snapshot is too old;
Listing~\ref{lst:app-pi-resend}):

% snippet: cpp/examples/pi/pi.cpp#resend [caption={Pi: chunks without an answer go out again},label={lst:app-pi-resend}]
\begin{lstlisting}[style=paglets-cpp,caption={Pi: chunks without an answer go out again \src{cpp/examples/pi/pi.cpp}},label={lst:app-pi-resend}]
// Chunks without an answer in time go out again (their host may be gone).
void tick() {
    if (worker_) return;
    const std::int64_t now = now_ms();
    for (auto& c : chunks_) {
        if (c.sent && !c.done && now - c.sent_ms > params_.chunk_timeout_ms) {
            c.sent = false;
            ++resent_;
        }
    }
    send_more();
    if (!finished()) (void)paglets::after_raw(500, "tick", {});
}
\end{lstlisting}
% /snippet

The conversion to decimal digits needs every hexadecimal digit, so it
runs once at the end, as passes of multiplying the fraction by $10^9$
whose carry is the next nine decimal digits, in slices of 200\,ms
(Listing~\ref{lst:pi-convert} in Section~\ref{sec:discussion-weak},
where the slicing is explained line by line).

% snippet: cpp/examples/pi/pi_msgs.hpp [caption={Pi's messages: the job's parameters, a worker's task, a result and the status},label={lst:app-pi-messages}]
\begin{lstlisting}[style=paglets-cpp,caption={Pi's messages: the job's parameters, a worker's task, a result and the status \src{cpp/examples/pi/pi\_msgs.hpp}},label={lst:app-pi-messages}]
// Messages of the pi compute example (planning/cpp-compute.md, section 5).

#pragma once

#include <cstdint>
#include <string>
#include <vector>

namespace pi_msgs {

// `start`: compute `digits` hex digits of pi in chunks of `chunk` digits,
// at most `max_in_flight` chunks at a time on hosts mesh-info selects. A
// chunk not answered within `chunk_timeout_ms` (its host may be gone) is
// sent again elsewhere. `work_ms` makes each chunk take at least that long
// (for demonstrations and tests).
//
// With `output_root`, the digits go, in order as they come in, to files
// below that named root of the job's host (a grant for `write` and
// `create`): `<output_path>.hex` and, with `decimal`, `<output_path>.txt`.
// The job then keeps only the chunks that came early, not all digits.
// With `decimal`, the job converts the hex digits into decimal digits at the
// end (about 1.2 per hex digit; the conversion needs every digit, so it
// cannot run per chunk).
struct Start {
    std::int64_t digits = 1000;
    std::int64_t chunk = 60;
    std::int64_t max_in_flight = 8;
    std::int64_t chunk_timeout_ms = 30'000;
    std::int64_t work_ms = 0;
    bool decimal = true;
    std::string output_root;  // empty: the digits stay in the job's memory
    std::string output_path = "pi";
};

struct Started {
    bool accepted = false;
    std::string reason;
};

// The arguments of a worker (a clone of the job paglet on another host).
struct Task {
    std::int64_t start = 0;
    std::int64_t count = 0;
    std::int64_t attempt = 0;
    std::int64_t work_ms = 0;
};

// `result`: a worker's digits.
struct Result {
    std::int64_t start = 0;
    std::int64_t count = 0;
    std::int64_t attempt = 0;
    std::string digits;
    std::string host;  // key ID of the host that computed it
    std::string host_name;
};

struct StatusRequest {};

// `status`: progress. `hex` and `decimal` show at most the first `shown`
// digits (the files have all of them).
struct Status {
    std::string state;        // idle, computing, converting, done, failed
    std::int64_t digits = 0;  // hex digits asked for
    std::int64_t done = 0;    // hex digits contiguous from the start
    std::int64_t chunks = 0;
    std::int64_t chunks_done = 0;
    std::int64_t in_flight = 0;
    std::int64_t resent = 0;  // chunks sent again (no answer in time)
    bool finished = false;    // every hex digit is there
    std::string hex;          // "3." followed by hex digits
    std::int64_t decimal_digits = 0;  // decimal digits that hex `digits` give
    std::int64_t decimal_done = 0;    // converted so far
    std::string decimal;              // "3." followed by decimal digits
    std::int64_t shown = 0;
    std::vector<std::string> files;  // the files written, below `output_root`
    std::vector<std::string> hosts;  // names of the hosts that computed chunks
    std::string error;
};

}  // namespace pi_msgs
\end{lstlisting}
% /snippet

A worker's result is a one-way \code{send} of a \code{Result} to the
endpoint it was given, followed by the release of its lease and
\code{dispose}. A late result for a chunk that is already done is
ignored, so a chunk that was merely slow is computed twice and the first
answer counts. The files are written through the \code{files} service
with a directory capability the job obtained from \code{grants} at the
start (Listing~\ref{lst:pi-grant}), dropped when the state reaches
\code{done}. Section~\ref{sec:evaluation} has the live run of this
example on a Mac, a Raspberry Pi~4 and a uConsole built on a Raspberry
Pi Compute Module~5.

\subsection{Mesh Chronicle}\label{sec:examples-chronicle}

The chronicle is a zombie \paglet{} (Section~\ref{sec:zombies}): it has a
replica on every host of the mesh, one replica samples the mesh and
writes down what changed, the others take over from the latest
checkpoint when it disappears, and when its lifetime is over the life log
goes to a sink host exactly once. The \paglet{}
(Listing~\ref{lst:app-chronicle}; its messages, work and journal entries
are Listing~\ref{lst:app-chronicle-messages}) derives from
\code{patterns::Zombie<Work>} and implements five hooks; the protocol,
the checkpoints, the takeovers and the delivery are in the base class.
Most of the file is \code{life\_log}, which renders the merged journal as
Markdown.

% snippet: cpp/examples/chronicle/chronicle.cpp [caption={The Mesh Chronicle},label={lst:app-chronicle},literate={–}{{--}}1]
\begin{lstlisting}[style=paglets-cpp,caption={The Mesh Chronicle \src{cpp/examples/chronicle/chronicle.cpp}},label={lst:app-chronicle},literate={–}{{--}}1]
// Mesh Chronicle (docs/examples/mesh-chronicle.md): a zombie paglet that
// writes the life log of its mesh. It has a replica on every host; the main
// samples mesh-info every few seconds and records what changed: hosts that
// disappear and come back, connections between hosts that are lost and
// restored (as each host reports them), and its own takeovers. Replicas
// that ran apart during a partition merge their journals when they meet, so
// what both sides saw ends up in one log. When its lifetime is over, the
// life log goes to the sink as Markdown, the owner is notified, and every
// replica ends.

#include "chronicle.schema.gen.hpp"

#include <paglets/patterns/notify.hpp>
#include <paglets/patterns/zombie.hpp>

#include <algorithm>
#include <map>
#include <numeric>
#include <set>
#include <string>
#include <vector>

using namespace paglets;
using namespace chronicle_msgs;
namespace zp = paglets::protocols::zombie;

namespace {

std::int64_t now_ms() {
    return paglets::now_ns() / 1'000'000;
}

std::string two(std::int64_t v) {
    return (v < 10 ? "0" : "") + std::to_string(v);
}

// "2026-10-13 14:05:09" (UTC).
std::string utc(std::int64_t ms) {
    std::int64_t secs = ms / 1000;
    std::int64_t days = secs / 86400;
    std::int64_t rest = secs % 86400;
    // Days to civil date (Howard Hinnant's algorithm).
    days += 719468;
    const std::int64_t era = (days >= 0 ? days : days - 146096) / 146097;
    const std::int64_t doe = days - era * 146097;
    const std::int64_t yoe = (doe - doe / 1460 + doe / 36524 - doe / 146096) / 365;
    std::int64_t y = yoe + era * 400;
    const std::int64_t doy = doe - (365 * yoe + yoe / 4 - yoe / 100);
    const std::int64_t mp = (5 * doy + 2) / 153;
    const std::int64_t d = doy - (153 * mp + 2) / 5 + 1;
    const std::int64_t m = mp < 10 ? mp + 3 : mp - 9;
    if (m <= 2) ++y;
    return std::to_string(y) + "-" + two(m) + "-" + two(d) + " " + two(rest / 3600) + ":" + two(rest % 3600 / 60) +
           ":" + two(rest % 60);
}

// "1 h 02 min", "3 min 05 s", "12 s".
std::string span(std::int64_t ms) {
    const std::int64_t s = std::max<std::int64_t>(0, ms) / 1000;
    if (s >= 3600) return std::to_string(s / 3600) + " h " + two(s % 3600 / 60) + " min";
    if (s >= 60) return std::to_string(s / 60) + " min " + two(s % 60) + " s";
    return std::to_string(s) + " s";
}

}  // namespace

class Chronicle final : public patterns::Zombie<Work> {
public:
    Chronicle() {
        router().on<Start>("start", [this](const Start& s, Message& m) -> std::int32_t {
            if (!group().empty()) return abi::bad_state;
            if (!s.sink.starts_with("host:") || s.lifetime_ms <= 0 || s.sample_ms < 1000) return abi::invalid_argument;
            Work w;
            w.sink = s.sink;
            w.lifetime_ms = s.lifetime_ms;
            w.sample_ms = s.sample_ms;
            w.started_ms = now_ms();
            patterns::ZombieTiming t;
            t.step_ms = s.sample_ms;
            if (s.beat_ms > 0) t.beat_ms = s.beat_ms;
            if (s.suspect_ms > 0) {
                t.suspect_ms = s.suspect_ms;
                t.jitter_ms = s.suspect_ms / 2;
            }
            if (s.round_ms > 0) t.round_ms = s.round_ms;
            if (s.expand_ms > 0) t.expand_ms = s.expand_ms;
            t.checkpoint_ms = std::max<std::int64_t>(s.sample_ms, t.beat_ms);
            start_group(s.sink, std::move(w), t);
            (void)m.reply(chronicle_msgs::Started{group(), {}});
            return abi::handled;
        });
        router().on("status", [this](Message& m) {
            const auto& r = replica();
            Status st;
            st.group = group();
            st.role = std::string(zp::to_string(r.role));
            st.replica = host_of(r.self);
            st.replicas = static_cast<std::uint32_t>(r.members.size());
            st.main = name_of(host_of(r.leader));
            st.samples = r.rank.progress;
            st.started_ms = work().started_ms;
            st.lifetime_ms = work().lifetime_ms;
            st.events = static_cast<std::int64_t>(r.journal.entries.size());
            (void)m.reply(st);
        });
        router().on("finish", [this](Message& m) {
            FinishReply f;
            f.accepted = acting();
            if (f.accepted) finish_ = true;
            const auto& r = replica();
            f.main = name_of(host_of(r.leader));
            (void)m.reply(f);
        });
    }

    void step(Work& w, std::function<void(bool done)> finished) override {
        auto mi = service("mesh-info");
        if (!mi) return finished(false);
        services::mesh_info::Client client{*mi};
        services::mesh_info::LandscapeRequest q;
        q.max_age_ms = std::max<std::int64_t>(3 * w.sample_ms, 3000);
        auto sent = client.landscape(q, [this, &w, finished](Result<services::mesh_info::Landscape> r, Message&) {
            if (!r) return finished(false);
            sample(w, r->hosts);
            const bool done = finish_ || now_ms() - w.started_ms >= w.lifetime_ms;
            if (done) note(Event{now_ms(), "end", here_, {}, std::to_string(w.samples) + " samples"});
            finished(done);
        });
        if (!sent) finished(false);
    }

    std::int64_t progress(const Work& w) const override { return w.samples; }

    std::string result_name() const override { return "chronicle-" + group() + ".md"; }

    void on_role(Role, Role to, Cause cause) override {
        if (to == Role::main) {
            takeover_ = true;
            takeover_cause_ = std::string(zp::to_string(cause));
        }
        if (to == Role::delivering) {
            patterns::notify(services::user_info::Level::info, "Mesh Chronicle " + group().substr(0, 8) + " is done",
                             "Its life log goes to " + work().sink + ".");
        }
    }

    Bytes result(const Work& w, const std::vector<zp::Entry>& journal) override {
        const std::string text = life_log(w, journal);
        return Bytes(text.begin(), text.end());
    }

private:
    void note(Event e) { (void)observe_value(e); }

    std::vector<Event> journal_events() const {
        std::vector<Event> out;
        for (const auto& [k, d] : replica().journal.entries) {
            Event e;
            if (decode(d, e)) out.push_back(std::move(e));
        }
        std::ranges::stable_sort(out, [](const Event& a, const Event& b) { return a.time_ms < b.time_ms; });
        return out;
    }

    std::string name_of(const std::string& key) const {
        const auto& w = work();
        for (std::size_t i = 0; i < w.members.size() && i < w.names.size(); ++i) {
            if (w.members[i] == key) return w.names[i];
        }
        return key.substr(0, 8);
    }

    // Compares the mesh now with the latest sample and records the changes.
    void sample(Work& w, const std::vector<services::mesh_info::Snapshot>& hosts) {
        const std::int64_t now = now_ms();
        if (!hosts.empty()) here_ = hosts.front().host;
        // Who is in the group: the replicas this main knows, with names from
        // mesh-info. Hosts that joined the mesh got a replica meanwhile.
        const bool first = w.members.empty();
        for (const auto& m : members()) {
            if (std::ranges::find(w.members, m) != w.members.end()) continue;
            w.members.push_back(m);
            w.names.emplace_back();
            if (!first) note(Event{now, "joined", m, here_});
        }
        for (const auto& h : hosts) {
            auto it = std::ranges::find(w.members, h.host);
            if (it != w.members.end() && !h.host_name.empty()) w.names[it - w.members.begin()] = h.host_name;
        }
        if (first) {
            std::string list;
            for (std::size_t i = 0; i < w.members.size(); ++i)
                list += (list.empty() ? "" : ", ") + (w.names[i].empty() ? w.members[i].substr(0, 8) : w.names[i]);
            note(Event{now, "begin", here_, {}, list});
        }
        if (takeover_) {
            takeover_ = false;
            if (w.samples > 0) {
                note(Event{now, "takeover", here_, {}, takeover_cause_, now - w.last_sample_ms, replica().reign});
            }
        }
        // What the merged journal says now: every lineage's observations,
        // the latest per host and per link.
        std::map<std::string, bool> host_up;
        std::set<std::string> down;  // "observer>other"
        for (const auto& m : w.members) host_up[m] = true;
        for (const auto& e : journal_events()) {
            if (e.kind == "host_down" || e.kind == "host_up") host_up[e.host] = e.kind == "host_up";
            if (e.kind == "link_down") down.insert(e.host + ">" + e.other);
            if (e.kind == "link_up") down.erase(e.host + ">" + e.other);
        }
        std::set<std::string> fresh;
        for (const auto& h : hosts) fresh.insert(h.host);
        for (const auto& [h, was_up] : host_up) {
            if (fresh.contains(h) && !was_up) note(Event{now, "host_up", h, here_});
            if (!fresh.contains(h) && was_up) note(Event{now, "host_down", h, here_});
        }
        // Links as each fresh host reports them (hosts that are not fresh
        // keep what they reported last).
        for (const auto& h : hosts) {
            for (const auto& l : h.links) {
                const std::string key = h.host + ">" + l.host;
                if (!l.up && !down.contains(key)) note(Event{now, "link_down", h.host, l.host, l.failure});
                if (l.up && down.contains(key)) note(Event{now, "link_up", h.host, l.host});
            }
        }
        if (w.samples > 0) w.alive_ms += std::min(now - w.last_sample_ms, 2 * w.sample_ms);
        w.last_sample_ms = now;
        ++w.samples;
    }

    // -- the life log --------------------------------------------------------------------------

    std::string life_log(const Work& w, const std::vector<zp::Entry>& journal) const {
        std::vector<Event> events;
        for (const auto& e : journal) {
            Event ev;
            if (decode(e.data, ev)) events.push_back(std::move(ev));
        }
        std::ranges::stable_sort(events, [](const Event& a, const Event& b) { return a.time_ms < b.time_ms; });
        // The same change reported by two lineages is one change.
        std::vector<Event> merged;
        for (auto& e : events) {
            const bool repeat = std::ranges::any_of(merged, [&](const Event& m) {
                if (e.kind == "joined") return m.kind == "joined" && m.host == e.host;  // a host joins once
                return m.kind == e.kind && m.host == e.host && m.other == e.other && e.time_ms - m.time_ms < 2 * w.sample_ms &&
                       e.kind != "takeover";
            });
            if (!repeat) merged.push_back(std::move(e));
        }
        auto name = [&](const std::string& key) {
            for (std::size_t i = 0; i < w.members.size() && i < w.names.size(); ++i) {
                if (w.members[i] == key && !w.names[i].empty()) return w.names[i];
            }
            return key.substr(0, 8);
        };
        const std::int64_t end = w.last_sample_ms;
        std::int64_t takeovers = 0, longest_gap = 0, dead = 0;
        for (const auto& e : merged) {
            if (e.kind != "takeover") continue;
            ++takeovers;
            longest_gap = std::max(longest_gap, e.gap_ms);
            dead += std::max<std::int64_t>(0, e.gap_ms - w.sample_ms);
        }
        std::string md;
        md += "# Mesh life log\n\n";
        md += "Written by the Mesh Chronicle, zombie group `" + group() + "`.\n\n";
        md += "| | |\n|---|---|\n";
        md += "| Began | " + utc(w.started_ms) + " UTC |\n";
        md += "| Ended | " + utc(end) + " UTC |\n";
        md += "| Lifetime | " + span(end - w.started_ms) + " |\n";
        std::string hosts;
        for (std::size_t i = 0; i < w.members.size(); ++i) hosts += (i ? ", " : "") + name(w.members[i]);
        md += "| Replicas | " + std::to_string(w.members.size()) + ": " + hosts + " |\n";
        md += "| Samples of the mesh | " + std::to_string(w.samples) + ", every " + span(w.sample_ms) + " |\n";
        md += "| Takeovers | " + std::to_string(takeovers) + " |\n";
        md += "| Time without a working main | " + span(dead) + " (longest gap " + span(longest_gap) + ") |\n\n";

        md += "## Timeline\n\n| Time (UTC) | What happened |\n|---|---|\n";
        for (const auto& e : merged) {
            std::string what;
            if (e.kind == "begin") {
                what = "The chronicle began on " + name(e.host) + ", with replicas on " + e.text + ".";
            } else if (e.kind == "host_down") {
                what = "**" + name(e.host) + "** disappeared (no fresh report seen from " + name(e.other) + ").";
            } else if (e.kind == "host_up") {
                what = "**" + name(e.host) + "** is back (seen from " + name(e.other) + ").";
            } else if (e.kind == "link_down") {
                what = name(e.host) + " lost its connection to " + name(e.other) +
                       (e.text.empty() ? "." : " (" + e.text + ").");
            } else if (e.kind == "link_up") {
                what = name(e.host) + " reaches " + name(e.other) + " again.";
            } else if (e.kind == "joined") {
                what = "**" + name(e.host) + "** joined the mesh and got a replica (seen from " + name(e.other) + ").";
            } else if (e.kind == "takeover") {
                what = "**" + name(e.host) + " took over as main** (" + e.text + ", reign " + std::to_string(e.term) +
                       ") after " + span(e.gap_ms) + " without a sample.";
            } else if (e.kind == "end") {
                what = "The work is done on " + name(e.host) + ": " + e.text + ".";
            } else {
                what = e.kind;
            }
            md += "| " + utc(e.time_ms) + " | " + what + " |\n";
        }

        // Connections: a pair is down while either end reports it down.
        md += "\n## Connections\n\n";
        struct Interval {
            std::int64_t from = 0, to = 0;
            std::string why;
        };
        std::map<std::string, std::vector<Interval>> pairs;
        std::map<std::string, std::set<std::string>> down_by;  // pair -> reporting directions
        auto pair_key = [&](const std::string& a, const std::string& b) {
            return name(std::min(a, b)) + " – " + name(std::max(a, b));
        };
        for (const auto& e : merged) {
            if (e.kind != "link_down" && e.kind != "link_up") continue;
            const auto key = pair_key(e.host, e.other);
            auto& by = down_by[key];
            const bool was = !by.empty();
            if (e.kind == "link_down") {
                by.insert(e.host);
            } else {
                by.erase(e.host);
            }
            if (!was && !by.empty()) pairs[key].push_back(Interval{e.time_ms, -1, e.text});
            if (was && by.empty() && !pairs[key].empty()) pairs[key].back().to = e.time_ms;
        }
        if (pairs.empty()) {
            md += "No connection was lost.\n";
        } else {
            md += "| Connection | Lost | Restored | For | Reason |\n|---|---|---|---|---|\n";
            for (const auto& [key, list] : pairs) {
                for (const auto& i : list) {
                    md += "| " + key + " | " + utc(i.from) + " | " + (i.to < 0 ? "not restored" : utc(i.to)) + " | " +
                          span((i.to < 0 ? end : i.to) - i.from) + " | " + (i.why.empty() ? "-" : i.why) + " |\n";
                }
            }
        }

        // Separations: hosts that could not see each other, from what either
        // side reported (a host gone from a main's view, a connection lost).
        // A side "worked" when a main there recorded something meanwhile; a
        // crashed host stays silent, a partition has work on both sides.
        md += "\n## Separations and partitions\n\n";
        std::set<std::pair<std::string, std::string>> cut;  // ordered pairs
        auto cut_pair = [&](const std::string& a, const std::string& b, bool on) {
            const auto p = std::pair{std::min(a, b), std::max(a, b)};
            if (on) {
                cut.insert(p);
            } else {
                cut.erase(p);
            }
        };
        auto components = [&]() {
            std::vector<std::string> hs = w.members;
            std::vector<std::size_t> parent(hs.size());
            std::iota(parent.begin(), parent.end(), 0);
            std::function<std::size_t(std::size_t)> find = [&](std::size_t x) {
                return parent[x] == x ? x : parent[x] = find(parent[x]);
            };
            for (std::size_t a = 0; a < hs.size(); ++a) {
                for (std::size_t b = a + 1; b < hs.size(); ++b) {
                    if (!cut.contains({std::min(hs[a], hs[b]), std::max(hs[a], hs[b])})) parent[find(a)] = find(b);
                }
            }
            std::map<std::size_t, std::vector<std::string>> groups;
            for (std::size_t i = 0; i < hs.size(); ++i) groups[find(i)].push_back(hs[i]);
            std::vector<std::vector<std::string>> out;
            for (auto& [root, g] : groups) out.push_back(std::move(g));
            return out;
        };
        struct Open {
            std::int64_t since = 0;
            std::vector<std::vector<std::string>> sides;
            std::set<std::string> active;  // observers that recorded something
        };
        std::optional<Open> open;
        std::vector<std::string> lines;
        std::int64_t partitions = 0;
        auto close = [&](std::int64_t at) {
            if (!open) return;
            if (at == open->since) {  // changes of one moment, not a separation
                open.reset();
                return;
            }
            std::string text;
            std::size_t worked = 0;
            for (const auto& side : open->sides) {
                std::string names;
                bool active = false;
                for (const auto& h : side) {
                    names += (names.empty() ? "" : ", ") + name(h);
                    active |= open->active.contains(h);
                }
                worked += active;
                text += (text.empty() ? "" : " / ") + std::string("{") + names + "} " + (active ? "worked" : "silent");
            }
            const bool partition = worked >= 2;
            partitions += partition;
            lines.push_back("| " + utc(open->since) + " | " + (at < 0 ? "not healed" : utc(at)) + " | " +
                            span((at < 0 ? end : at) - open->since) + " | " + text + " | " +
                            (partition ? "partition: both sides had a main" : "a host was gone or cut off") + " |");
            open.reset();
        };
        for (const auto& e : merged) {
            if (e.kind == "host_down") cut_pair(e.host, e.other, true);
            if (e.kind == "host_up") cut_pair(e.host, e.other, false);
            if (e.kind == "link_down") cut_pair(e.host, e.other, true);
            if (e.kind == "link_up") cut_pair(e.host, e.other, false);
            const std::string observer = e.kind == "host_down" || e.kind == "host_up" ? e.other : e.host;
            const auto parts = components();
            if (parts.size() < 2) {
                close(e.time_ms);
                continue;
            }
            if (!open || open->sides != parts) {
                close(e.time_ms);
                open = Open{e.time_ms, parts, {}};
            }
            open->active.insert(observer);
        }
        close(-1);
        if (lines.empty()) {
            md += "The hosts could always see each other.\n";
        } else {
            md += "| From | To | For | Sides | What it was |\n|---|---|---|---|---|\n";
            for (const auto& l : lines) md += l + "\n";
            md += "\n" + std::to_string(partitions) +
                  " partition(s) with a main on both sides. What both sides observed is merged into this log: the "
                  "replicas that worked apart handed over their journals when they met again; the lineage that got "
                  "furthest went on, the others forgot their work.\n";
        }
        return md;
    }

    std::string here_;
    bool takeover_ = false;
    std::string takeover_cause_;
    bool finish_ = false;
};

PAGLETS_PAGLET(Chronicle)
\end{lstlisting}
% /snippet

% snippet: cpp/examples/chronicle/chronicle_msgs.hpp [caption={The chronicle's messages, work and journal entries},label={lst:app-chronicle-messages}]
\begin{lstlisting}[style=paglets-cpp,caption={The chronicle's messages, work and journal entries \src{cpp/examples/chronicle/chronicle\_msgs.hpp}},label={lst:app-chronicle-messages}]
// Messages of the Mesh Chronicle (docs/examples/mesh-chronicle.md): a zombie
// paglet that writes the life log of its mesh.

#pragma once

#include <cstdint>
#include <string>
#include <vector>

namespace chronicle_msgs {

// `start`: begin the chronicle. It lives `lifetime_ms`, samples the mesh
// every `sample_ms` and delivers its life log to `sink` (`host:NAME`). The
// timing of the protocol can be shortened for tests (0: the defaults).
struct Start {
    std::string sink;
    std::int64_t lifetime_ms = 24LL * 3600 * 1000;
    std::int64_t sample_ms = 10'000;
    std::int64_t beat_ms = 0;
    std::int64_t suspect_ms = 0;
    std::int64_t round_ms = 0;
    std::int64_t expand_ms = 0;  // how often the main looks for hosts without a replica
};

struct Started {
    std::string group;
    std::vector<std::string> hosts;  // where the replicas are
};

// The state of the work (replicated with every checkpoint).
struct Work {
    std::string sink;
    std::int64_t lifetime_ms = 0;
    std::int64_t sample_ms = 0;
    std::int64_t started_ms = 0;      // when the group was created
    std::int64_t samples = 0;         // progress
    std::int64_t last_sample_ms = 0;  // the latest sample of the lineage
    std::int64_t alive_ms = 0;        // time with a working main, summed over samples
    std::vector<std::string> members;     // host key IDs
    std::vector<std::string> names;       // their names, same order
};

// One entry of the journal: something the main observed.
struct Event {
    std::int64_t time_ms = 0;
    std::string kind;   // begin, joined, host_up, host_down, link_down, link_up, takeover, end
    std::string host;   // key ID: the host it is about (takeover: the new main's host)
    std::string other;  // links: the other end; hosts: the main that saw it
    std::string text;   // details: a failure, the cause of a takeover
    std::int64_t gap_ms = 0;  // takeover: time since the lineage's latest sample
    std::uint64_t term = 0;   // takeover: the new reign
};

struct StatusRequest {};

struct Status {
    std::string group;
    std::string role;   // of this replica
    std::string replica;  // this replica's host (key ID)
    std::uint32_t replicas = 0;  // replicas it knows of
    std::string main;   // host name of the main this replica follows (itself if main)
    std::int64_t samples = 0;
    std::int64_t started_ms = 0;
    std::int64_t lifetime_ms = 0;
    std::int64_t events = 0;  // entries in this replica's journal
};

struct FinishRequest {};

struct FinishReply {
    bool accepted = false;  // only the acting main finishes; others name it
    std::string main;
};

}  // namespace chronicle_msgs
\end{lstlisting}
% /snippet

The \code{start} handler returns a status code instead of a reply when
the group exists or the arguments are wrong (\code{bad\_state},
\code{invalid\_argument}); otherwise it fills the \code{Work}, the state
that every checkpoint replicates, shortens the protocol's timing when
asked (the tests do), and calls \code{start\_group(sink, work, timing)}.
The group starts after \code{mesh-info} answered, so the reply knows the
group ID but not the hosts; \code{Started.hosts} stays empty.
\code{status} reads the protocol state through \code{replica()}: the
role, the members, the main it follows, the rank's progress and the
journal.
\code{finish} is accepted on the acting main only.

The base class calls \code{step} on the acting main, with a pause of
\code{step\_ms} between steps, which the chronicle sets to the sample
interval (\code{t.step\_ms = s.sample\_ms}). A step may wait for a reply, as this one
waits for the landscape; if the main stopped meanwhile, its
\code{finished} call is ignored. \code{sample} compares the fresh
landscape with the view derived from \code{journal\_events()}, the merged
journal of every lineage, not with the work of the current lineage, so
that after a partition the winner sees what the other side recorded.
Every observation is one \code{Event} written with \code{observe\_value},
which the base class accepts on the acting main only. \code{progress} is
the number of samples; when lineages meet, a finished one wins, then the
one with more progress, then the later reign (Section~\ref{sec:zombies}).
\code{on\_role} notes a takeover for the next sample and notifies the
owner when the replica delivers. \code{result} builds the life log from
the final work and the merged journal: a change that two lineages both
recorded within two sample intervals is one event, lost connections
become intervals per pair of hosts, and the separations are the
connected components of the hosts under the recorded cuts, each side
marked as worked or silent, which tells a partition (a main on both
sides) from a crash. The \code{Work} is small on purpose, names and
counters, because the winner's work replaces the loser's after a
partition; what must not be lost goes into the journal.
